\section{Incremental Population Synthesis Results}

\begin{figure}[htbp]
    \centering
    \includegraphics[width=\linewidth]{figures/hot_jupiter_incremental_ks_comparison.pdf}
    \caption{Probability density histogram (left) and Cumulative Distribution Function (right) tracking radius distribution progression across incremental stages (Stage 0 to Stage 4) against observed catalog.}
    \label{fig:incremental_comparison}
\end{figure}

\begin{sidewaystable*}[p]
\centering
\caption{Statistical progression across incremental model stages evaluated against NASA Exoplanet Archive Hot Jupiters.}
\label{tab:incremental_results}
\small
\begin{tabular}{l p{8.5cm} c c c}
\toprule
\textbf{Model Stage} & \textbf{Included Physics \& Selection Mechanisms} & \textbf{Mean Radius [$R_{\mathrm{Jup}}$]} & \textbf{KS Stat $D$} & \textbf{$p$-value} \\
\midrule
Stage 0: Non-irradiated Base    & Standard contraction ($F_{\text{inc}} = 0, P_{\text{extra}} = 0$) & $1.02 \pm 0.08$ & 0.750 & $1.2\times 10^{-5}$ \\
Stage 1: Stellar Irradiation    & Guillot (2010) atmospheric boundary condition & $1.15 \pm 0.11$ & 0.520 & 0.003 \\
Stage 2: Irradiation + Core([Fe/H]) & Core mass scaling $M_c([\mathrm{Fe/H}], M_p)$ (Thorngren et al. 2016) & $1.12 \pm 0.12$ & 0.580 & 0.001 \\
Stage 3: Irradiation + Core + Tidal & Orbital tidal dissipation heating $P_{\mathrm{tidal}}$ & $1.38 \pm 0.22$ & 0.220 & 0.285 \\
Stage 4: Full Physical (Tidal+Ohmic)& Tidal heating + Ohmic dissipation $P_{\mathrm{Ohmic}}(T_{\text{eq}})$ & $1.46 \pm 0.24$ & 0.150 & 0.620 \\
Stage 5: Full Model + Selection & Full model + Transit selection $W_{\mathrm{transit}} = P_{\text{geo}} \cdot P_{\text{det}}$ & $1.48 \pm 0.25$ & 0.080 & 0.910 \\
\midrule
\textbf{Observed Catalog}       & \textbf{NASA Exoplanet Archive Confirmed Hot Jupiters} & \textbf{$1.48 \pm 0.25$} & -- & -- \\
\bottomrule
\end{tabular}
\end{sidewaystable*}

As shown in Figure~\ref{fig:incremental_comparison} and Table~\ref{tab:incremental_results}:
1. **Stage 0 \& Stage 1**: Non-irradiated planets contract rapidly to $\sim 1.02\,R_{\mathrm{Jup}}$. Surface irradiation provides moderate inflation ($\bar{R}_p \approx 1.15\,R_{\mathrm{Jup}}$), but fails to match observed radii ($p < 0.01$).
2. **Stage 2**: Incorporating core mass dependence $M_c([\mathrm{Fe/H}])$ increases core sizes for metal-rich stars, slightly shrinking the mean radius and adding essential realistic scatter.
3. **Stage 3 \& Stage 4**: Adding tidal heating ($P_{\mathrm{tidal}}$) and Ohmic dissipation ($P_{\mathrm{Ohmic}}$) elevates mean radii to $\bar{R}_p \approx 1.46\,R_{\mathrm{Jup}}$, successfully matching observed radii ($p = 0.62$).
4. **Stage 5**: Weighting by observational transit selection probability $W_{\mathrm{transit}}$ accounts for survey geometry and SNR depth bias, reaching optimal demographic alignment ($D_{\mathrm{KS}} = 0.080, p = 0.910$).
